% =============================================================================
% MusicTUI: A Terminal-Based Music Player with Synced Lyrics and ANSI Art
% -----------------------------------------------------------------------------
% Restyled edition.  Content, structure, code listings,
% project macros and bibliography are unchanged; only the presentation layer
% (preamble, environments, spacing, colour) has been redesigned.
%
% Build:  pdflatex musictui_paper.tex   (run twice for references)
% Needs only standard TeX Live packages.  Optional packages (fontawesome5,
% inconsolata/zi4, microtype) are detected automatically and replaced with
% safe fall-backs when absent.
% =============================================================================
\documentclass[a4paper,11pt,twoside]{article}

%% ---------------------------------------------------------------------------
%% 1.  Encoding, fonts, microtypography
%% ---------------------------------------------------------------------------
\usepackage[T1]{fontenc}
\usepackage[utf8]{inputenc}
\usepackage{amsmath,amssymb}
\usepackage[sc]{mathpazo}                 % Palatino text + math (serif body)
\usepackage[scaled=0.92]{helvet}          % Helvetica for headings/captions
\IfFileExists{zi4.sty}{%
  \usepackage[scaled=0.93]{zi4}%          % Inconsolata (preferred monospace)
}{\IfFileExists{inconsolata.sty}{%
  \usepackage[scaled=0.93]{inconsolata}%
}{%
  % universally available fall-back: Courier (URW) scaled down to Palatino's colour
  \DeclareFontFamily{T1}{pcs}{\hyphenchar\font=-1 }%
  \DeclareFontShape{T1}{pcs}{m}{n}{<->s*[0.86]pcrr8t}{}%
  \DeclareFontShape{T1}{pcs}{m}{sl}{<->s*[0.86]pcrro8t}{}%
  \DeclareFontShape{T1}{pcs}{m}{it}{<->ssub*pcs/m/sl}{}%
  \DeclareFontShape{T1}{pcs}{b}{n}{<->s*[0.86]pcrb8t}{}%
  \DeclareFontShape{T1}{pcs}{b}{sl}{<->s*[0.86]pcrbo8t}{}%
  \DeclareFontShape{T1}{pcs}{b}{it}{<->ssub*pcs/b/sl}{}%
  \DeclareFontShape{T1}{pcs}{bx}{n}{<->ssub*pcs/b/n}{}%
  \DeclareFontShape{T1}{pcs}{bx}{sl}{<->ssub*pcs/b/sl}{}%
  \DeclareFontShape{T1}{pcs}{bx}{it}{<->ssub*pcs/b/sl}{}%
  \renewcommand{\ttdefault}{pcs}%
}}
\IfFileExists{microtype.sty}{%
  \usepackage[protrusion=true,expansion=true,final]{microtype}%
}{}

%% ---------------------------------------------------------------------------
%% 2.  Colour palette  (deep blue + teal + warm amber, plus greys)
%% ---------------------------------------------------------------------------
\usepackage[table]{xcolor}
\definecolor{accent}{HTML}{14507A}        % primary: deep blue
\definecolor{accentb}{HTML}{0E7C7B}       % secondary: teal
\definecolor{warm}{HTML}{A8611A}          % tertiary: muted amber
\definecolor{ink}{HTML}{1B2430}           % near-black for headings
\definecolor{codebg}{HTML}{F5F7F9}        % code block background
\definecolor{codeframe}{HTML}{D3DAE1}     % code block frame / rules
\definecolor{codegray}{HTML}{6A737D}      % comments, line numbers

%% ---------------------------------------------------------------------------
%% 3.  Core packages
%% ---------------------------------------------------------------------------
\usepackage{graphicx}
\usepackage[a4paper,left=1.15in,right=1.15in,top=1in,bottom=1in,
            headheight=15pt,headsep=0.3in,footskip=0.4in]{geometry}
\usepackage{setspace}
\usepackage{parskip}
\usepackage{etoolbox}
\usepackage{array}
\usepackage{tabularx}
\usepackage{booktabs}
\usepackage{float}
\usepackage{caption}
\usepackage{subcaption}
\usepackage{enumitem}
\usepackage{tikz}
\usepackage{pgfplots}
\pgfplotsset{compat=newest}
\usepackage[ruled]{algorithm}
\usepackage{algpseudocode}
\usepackage{listings}
\usepackage{listingsutf8}
\usepackage[skins,breakable]{tcolorbox}
\usepackage{titlesec}
\usepackage{fancyhdr}

%% Icons for the call-out boxes: Font Awesome when installed, Dingbats otherwise
\IfFileExists{fontawesome5.sty}{%
  \usepackage{fontawesome5}%
  \newcommand{\iconDesign}{\faLightbulb}%
  \newcommand{\iconArch}{\faSitemap}%
  \newcommand{\iconTest}{\faFlask}%
}{%
  \usepackage{pifont}%
  \newcommand{\iconDesign}{\ding{72}}%
  \newcommand{\iconArch}{\ding{117}}%
  \newcommand{\iconTest}{\ding{51}}%
}

\usepackage{hyperref}                     % keep near the end
\hypersetup{
  colorlinks=true,
  linkcolor=accent, citecolor=accentb, urlcolor=accent,
  bookmarksnumbered=true, pdfpagemode=UseOutlines,
  pdftitle={MusicTUI: A Terminal-Based Music Player with Synced Lyrics, ANSI Half-Block Colour Art, and Event-Driven Architecture},
  pdfauthor={MusicTUI Project},
  pdfkeywords={terminal user interface, music player, Python, Textual, synced lyrics, ANSI art, event-driven architecture, dependency injection, SQLite, VLC}
}

%% ---------------------------------------------------------------------------
%% 4.  Page layout, spacing, paragraphs
%% ---------------------------------------------------------------------------
\onehalfspacing
\setlength{\parskip}{6pt plus 1.5pt minus 1pt}
\setlength{\parindent}{0pt}
\raggedbottom
\setlength{\emergencystretch}{2em}
\widowpenalty=10000
\clubpenalty=10000
\setlength{\floatsep}{14pt plus 3pt minus 3pt}
\setlength{\textfloatsep}{16pt plus 3pt minus 4pt}
\setlength{\intextsep}{14pt plus 3pt minus 3pt}

%% ---------------------------------------------------------------------------
%% 5.  Section headings (titlesec)
%% ---------------------------------------------------------------------------
\newcommand{\secbar}{\textcolor{accent}{\rule[-0.32em]{4pt}{1.5em}}\hspace{0.6em}}

\titleformat{\section}[block]
  {\normalfont\sffamily\Large\bfseries\color{ink}}
  {\secbar\textcolor{accent}{\thesection}}{0.6em}{}
  [\vspace{-0.15em}{\color{accent!25}\titlerule[0.6pt]}]
\titleformat{name=\section,numberless}[block]
  {\normalfont\sffamily\Large\bfseries\color{ink}}
  {}{0pt}{\secbar}
  [\vspace{-0.15em}{\color{accent!25}\titlerule[0.6pt]}]
\titlespacing*{\section}{0pt}{3.4ex plus 1ex minus .4ex}{1.8ex plus .3ex}

\titleformat{\subsection}[block]
  {\normalfont\sffamily\large\bfseries\color{ink}}
  {\textcolor{accentb}{\thesubsection}}{0.55em}{}
\titlespacing*{\subsection}{0pt}{2.8ex plus .8ex minus .3ex}{1.1ex plus .2ex}

\titleformat{\subsubsection}[block]
  {\normalfont\sffamily\normalsize\bfseries\color{ink!85}}
  {\textcolor{accentb!85!black}{\thesubsubsection}}{0.5em}{}
\titlespacing*{\subsubsection}{0pt}{2.2ex plus .6ex minus .2ex}{0.8ex plus .2ex}

\titleformat{\paragraph}[runin]
  {\normalfont\sffamily\normalsize\bfseries\color{accent}}
  {}{0pt}{}
\titlespacing*{\paragraph}{0pt}{1.6ex plus .5ex minus .2ex}{0.8em}

%% ---------------------------------------------------------------------------
%% 6.  Running headers / footers
%%     odd pages : current section (left)  | page number (right)
%%     even pages: page number (left)      | short title (right)
%% ---------------------------------------------------------------------------
\pagestyle{fancy}
\fancyhf{}
\renewcommand{\sectionmark}[1]{\markboth{\thesection\enspace #1}{\thesection\enspace #1}}
\fancyhead[LO]{\sffamily\small\color{codegray}\leftmark}
\fancyhead[RO]{\sffamily\small\bfseries\color{accent}\thepage}
\fancyhead[LE]{\sffamily\small\bfseries\color{accent}\thepage}
\fancyhead[RE]{\sffamily\small\color{codegray}MusicTUI}
\renewcommand{\headrulewidth}{0.4pt}
\renewcommand{\headrule}{{\color{accent!35}\hrule width\headwidth height\headrulewidth}}
\renewcommand{\footrulewidth}{0pt}
\fancypagestyle{plain}{%
  \fancyhf{}%
  \fancyfoot[C]{\sffamily\small\color{codegray}\thepage}%
  \renewcommand{\headrulewidth}{0pt}%
}

%% ---------------------------------------------------------------------------
%% 7.  Captions
%% ---------------------------------------------------------------------------
\captionsetup{
  font=small, labelfont={sf,bf,color=accent}, textfont={},
  labelsep=period, justification=justified, singlelinecheck=true,
  margin=0.4em
}
\captionsetup[table]{position=above, skip=7pt}
\captionsetup[figure]{position=below, skip=9pt}
\captionsetup[algorithm]{position=above}
\captionsetup[ruled]{font=small, labelfont={sf,bf,color=accent}, labelsep=period,
  singlelinecheck=false, justification=raggedright}
\captionsetup[lstlisting]{
  font=small, labelfont={sf,bf,color=accent}, labelsep=period,
  justification=raggedright, singlelinecheck=false, skip=4pt, margin=0pt
}

%% ---------------------------------------------------------------------------
%% 8.  Lists
%% ---------------------------------------------------------------------------
\setlist{itemsep=0.2em plus 0.05em, parsep=0pt, partopsep=0pt,
         topsep=0.35em plus 0.15em minus 0.1em}
\setlist[itemize,1]{leftmargin=1.7em, labelsep=0.55em,
  label=\raisebox{0.18ex}{\textcolor{accent}{\scriptsize$\blacksquare$}}}
\setlist[itemize,2]{leftmargin=1.5em, labelsep=0.5em,
  label=\textcolor{accentb}{\textbf{--}}}
\setlist[itemize,3]{leftmargin=1.4em, label=\textcolor{codegray}{$\circ$}}
\setlist[enumerate,1]{leftmargin=2em, labelsep=0.55em,
  label=\textcolor{accent}{\sffamily\bfseries\arabic*.}}
\setlist[enumerate,2]{leftmargin=2em, label=\textcolor{accentb}{\alph*)}}

%% ---------------------------------------------------------------------------
%% 9.  Tables  (booktabs only: no vertical rules, no \hline)
%% ---------------------------------------------------------------------------
\setlength{\heavyrulewidth}{1.1pt}
\setlength{\lightrulewidth}{0.5pt}
\setlength{\aboverulesep}{0.45ex}
\setlength{\belowrulesep}{0.65ex}
\AtBeginEnvironment{table}{%
  \small\singlespacing
  \renewcommand{\arraystretch}{1.3}%
  \setlength{\tabcolsep}{8pt}%
  \rowcolors{2}{white}{accent!5}%
}

%% ---------------------------------------------------------------------------
%% 10.  Algorithms
%% ---------------------------------------------------------------------------
\floatname{algorithm}{Algorithm}
\AtBeginEnvironment{algorithmic}{\small\singlespacing}
\algrenewcommand{\alglinenumber}[1]{\sffamily\scriptsize\color{codegray}#1:}
\algrenewcommand{\algorithmicfunction}{\textcolor{accent}{\textbf{function}}}
\algrenewcommand{\algorithmicif}{\textcolor{accent}{\textbf{if}}}
\algrenewcommand{\algorithmicthen}{\textcolor{accent}{\textbf{then}}}
\algrenewcommand{\algorithmicend}{\textcolor{accent}{\textbf{end}}}
\algrenewcommand{\algorithmicreturn}{\textcolor{accent}{\textbf{return}}}

%% ---------------------------------------------------------------------------
%% 11.  Code listings  (listings + tcolorbox, no shell-escape needed)
%% ---------------------------------------------------------------------------
\newcommand{\codefont}{\ttfamily\footnotesize\renewcommand{\baselinestretch}{1.12}\selectfont}

\lstdefinelanguage{Python}{
  morekeywords={and,as,assert,async,await,break,case,class,continue,def,del,
    elif,else,except,finally,for,from,global,if,import,in,is,lambda,match,
    None,nonlocal,not,or,pass,raise,return,True,False,try,while,with,yield},
  morekeywords=[2]{int,str,bool,float,dict,list,set,tuple,object,bytes,Any,
    Self,Enum,Message,print,len,isinstance,hasattr,getattr,setattr,round,max,
    min,zip,range,super,type,hash,property,staticmethod,classmethod,
    Exception,ValueError,ImportError,NotImplemented},
  morekeywords=[3]{self,cls},
  sensitive=true,
  morecomment=[l]\#,
  morestring=[b]",
  morestring=[b]',
  moredelim=[l][\color{accentb}]{@},
}
\lstdefinelanguage{Console}{
  morecomment=[l]\#,
  literate={\$\ }{{\textcolor{accentb}{\$\ }}}2,
}
\lstset{
  basicstyle=\codefont,
  keywordstyle=\color{accent}\bfseries,
  keywordstyle=[2]\color{accentb},
  keywordstyle=[3]\color{codegray},
  commentstyle=\color{codegray},
  stringstyle=\color{warm},
  identifierstyle=\color{ink},
  showstringspaces=false,
  breaklines=true, breakatwhitespace=false,
  postbreak=\mbox{\textcolor{accentb}{\tiny$\hookrightarrow$}\space},
  breakindent=0pt, breakautoindent=true,
  tabsize=4, columns=fullflexible, keepspaces=true,
  captionpos=t, aboveskip=0pt, belowskip=0pt,
  frame=l, framerule=0.5pt, framesep=5pt, rulecolor=\color{codeframe},
  numbers=left, stepnumber=1, numbersep=10pt,
  numberstyle=\tiny\ttfamily\color{codegray},
  xleftmargin=22pt, xrightmargin=0pt,
}
\lstdefinestyle{console}{language=Console, numbers=none, frame=none,
  xleftmargin=0pt, breaklines=true}

%% Shared look of code blocks (also used for the two `verbatim` trees)
\tcbset{codebox/.style={
  enhanced, breakable, colback=codebg, colframe=codeframe, boxrule=0.5pt,
  arc=3pt, sharp corners=west, borderline west={2.6pt}{0pt}{accent},
  left=5pt, right=6pt, top=5pt, bottom=5pt,
  before skip=0.7em plus 2pt, after skip=1.0em plus 2pt,
  pad at break*=2mm, toprule at break=0pt, bottomrule at break=0pt,
}}
\BeforeBeginEnvironment{lstlisting}{\begin{tcolorbox}[codebox]}
\AfterEndEnvironment{lstlisting}{\end{tcolorbox}}

\BeforeBeginEnvironment{verbatim}{\begin{tcolorbox}[codebox]}
\AfterEndEnvironment{verbatim}{\end{tcolorbox}}
\AtBeginEnvironment{verbatim}{\codefont\parskip=0pt\relax}

%% Box-drawing glyphs used by the two directory/wiring trees (drawn with rules
%% so they connect perfectly in any monospace font and need no special font).
\makeatletter
\newdimen\bd@w \newdimen\bd@y \newdimen\bd@t
\newcommand{\bd@setup}{%
  \leavevmode
  \bd@w=\fontcharwd\font`0\relax
  \bd@y=0.5\fontdimen5\font\relax
  \bd@t=0.45pt\relax
}
\newcommand{\bd@vfull}{\vrule height\ht\strutbox depth\dp\strutbox width\bd@t}
\newcommand{\bd@place}[2]{\rlap{\kern#1\relax#2}}
\newcommand{\bdVert}{\bd@setup\bd@place{\dimexpr.5\bd@w-.5\bd@t\relax}{\bd@vfull}\kern\bd@w}
\newcommand{\bdHorz}{\bd@setup\bd@place{0pt}{\raise\bd@y\hbox{\vrule height\bd@t width\bd@w}}\kern\bd@w}
\newcommand{\bdTee}{\bd@setup
  \bd@place{\dimexpr.5\bd@w-.5\bd@t\relax}{\bd@vfull}%
  \bd@place{.5\bd@w}{\raise\bd@y\hbox{\vrule height\bd@t width.5\bd@w}}\kern\bd@w}
\newcommand{\bdElbow}{\bd@setup
  \bd@place{\dimexpr.5\bd@w-.5\bd@t\relax}{\raise\bd@y\hbox{\vrule height\dimexpr\ht\strutbox-\bd@y\relax depth0pt width\bd@t}}%
  \bd@place{.5\bd@w}{\raise\bd@y\hbox{\vrule height\bd@t width.5\bd@w}}\kern\bd@w}
\newcommand{\bdTopRight}{\bd@setup
  \bd@place{0pt}{\raise\bd@y\hbox{\vrule height\bd@t width.5\bd@w}}%
  \bd@place{\dimexpr.5\bd@w-.5\bd@t\relax}{\lower\dp\strutbox\hbox{\vrule height\dimexpr\dp\strutbox+\bd@y\relax depth0pt width\bd@t}}\kern\bd@w}
\newcommand{\bdBotRight}{\bd@setup
  \bd@place{0pt}{\raise\bd@y\hbox{\vrule height\bd@t width.5\bd@w}}%
  \bd@place{\dimexpr.5\bd@w-.5\bd@t\relax}{\raise\bd@y\hbox{\vrule height\dimexpr\ht\strutbox-\bd@y\relax depth0pt width\bd@t}}\kern\bd@w}
\makeatother
\DeclareUnicodeCharacter{2502}{\bdVert}
\DeclareUnicodeCharacter{2500}{\bdHorz}
\DeclareUnicodeCharacter{251C}{\bdTee}
\DeclareUnicodeCharacter{2514}{\bdElbow}
\DeclareUnicodeCharacter{2510}{\bdTopRight}
\DeclareUnicodeCharacter{2518}{\bdBotRight}

%% ---------------------------------------------------------------------------
%% 12.  Call-out boxes
%% ---------------------------------------------------------------------------
\tcbset{callout/.style={
  enhanced jigsaw, breakable, arc=3pt, boxrule=0.6pt,
  colframe=#1, colback=#1!5!white, colbacktitle=#1, coltitle=white,
  fonttitle=\sffamily\bfseries\small, toptitle=2.2pt, bottomtitle=2.2pt,
  left=9pt, right=9pt, top=6pt, bottom=6pt,
  before skip=0.9em plus 2pt, after skip=1.1em plus 2pt,
  drop fuzzy shadow=black!18,
  before upper={\raggedright\sloppy\setstretch{1.12}\setlength{\parskip}{4pt}\vspace*{-\parskip}},
}}
\newtcolorbox{designbox}{callout=accent,
  title={\iconDesign\hspace{0.55em}Design Decision}}
\newtcolorbox{archbox}{callout=accentb,
  title={\iconArch\hspace{0.55em}Architecture Note}}
\newtcolorbox{testbox}{callout=warm,
  title={\iconTest\hspace{0.55em}Testing Note}}

%% ---------------------------------------------------------------------------
%% 13.  Figure placeholder helper  (\figfbox wraps placeholder content in a styled box)
%% ---------------------------------------------------------------------------
\newcommand{\figfbox}[1]{%
  \tcbox[enhanced, nobeforeafter, colback=accent!4, colframe=accent!45,
    boxrule=0pt, arc=3pt, boxsep=6pt, left=2pt, right=2pt, top=2pt, bottom=2pt,
    fontupper=\color{black!70},
    frame code={\path[draw=accent!50, dash pattern=on 4pt off 3pt,
      line width=0.7pt, rounded corners=3pt]
      (frame.south west) rectangle (frame.north east);}]{#1}}
\AtBeginEnvironment{figure}{\let\fbox\figfbox\singlespacing}

%% ---------------------------------------------------------------------------
%% 14.  Abstract + keywords
%% ---------------------------------------------------------------------------
\renewtcolorbox{abstract}{
  enhanced, breakable, colback=accent!5, colframe=accent!5, boxrule=0pt,
  arc=2pt, borderline west={3pt}{0pt}{accent},
  left=11pt, right=11pt, top=8pt, bottom=8pt,
  before skip=1.4em, after skip=0.6em,
  before upper={\singlespacing\small\setlength{\parskip}{3pt}%
    \noindent{\sffamily\bfseries\color{accent}\textls[80]{ABSTRACT}}\par\vspace{0.3em}},
}
\providecommand{\textls}[2][]{#2}   % letter-spacing fall-back when microtype is absent
\newcommand{\keywords}[1]{%
  \par\vspace{0.2em}\noindent
  {\singlespacing\small
   \textcolor{accent}{\sffamily\bfseries Keywords}\hspace{0.8em}%
   \textcolor{black!80}{#1}\par}%
  \vspace{0.2em}}

%% ---------------------------------------------------------------------------
%% 15.  Title block
%% ---------------------------------------------------------------------------
\makeatletter
\renewcommand{\maketitle}{%
  \begingroup
  \thispagestyle{plain}%
  \singlespacing
  \par\noindent{\color{accent}\rule{\linewidth}{2.6pt}}\par
  \vspace{1.3em}%
  {\raggedright\sffamily\bfseries\fontsize{15}{20}\selectfont\color{accent!85!black}\@title\par}%
  \vspace{0.9em}%
  {\color{accentb}\rule{2.6cm}{1.5pt}}\par
  \vspace{0.9em}%
  {\raggedright\sffamily\large\color{ink}\@author\par}%
  \vspace{0.5em}%
  {\raggedright\sffamily\small\color{codegray}\@date\par}%
  \vspace{0.8em}%
  \noindent{\color{accent!30}\rule{\linewidth}{0.5pt}}\par
  \endgroup
}
\makeatother

%% ---------------------------------------------------------------------------
%% 16.  Bibliography
%% ---------------------------------------------------------------------------
\makeatletter
\renewcommand{\@biblabel}[1]{\textcolor{accent}{\sffamily\bfseries[#1]}}
\makeatother
\AtBeginEnvironment{thebibliography}{\small\singlespacing\setlength{\itemsep}{3pt plus 1pt}}
\apptocmd{\thebibliography}{\markboth{References}{References}}{}{}

%% ---------------------------------------------------------------------------
%% 17.  Project-specific macros (UNCHANGED from the original source)
%% ---------------------------------------------------------------------------
\newcommand{\py}{\texttt{Python}}
\newcommand{\textual}{\texttt{Textual}}
\newcommand{\vlc}{\texttt{VLC}}
\newcommand{\sqlite}{\texttt{SQLite}}
\newcommand{\lrc}{\texttt{LRC}}
\newcommand{\ansi}{\texttt{ANSI}}
\newcommand{\ytmusic}{\texttt{ytmusicapi}}
\newcommand{\ytdlp}{\texttt{yt-dlp}}
\newcommand{\pil}{\texttt{PIL}}
\newcommand{\tui}{\texttt{TUI}}
\newcommand{\musictui}{\texttt{MusicTUI}}
\newcommand{\eventbus}{\texttt{EventBus}}
\newcommand{\protocol}{\texttt{Protocol}}
\newcommand{\track}{\texttt{Track}}
\newcommand{\lyric}{\texttt{LyricLine}}
\newcommand{\playlist}{\texttt{Playlist}}
\newcommand{\playstate}{\texttt{PlayState}}
\newcommand{\repeatmode}{\texttt{RepeatMode}}
\newcommand{\audiobackend}{\texttt{AudioBackend}}
\newcommand{\catalogsource}{\texttt{CatalogSource}}
\newcommand{\trackcache}{\texttt{TrackCache}}
\newcommand{\musicqueue}{\texttt{MusicQueue}}
\newcommand{\player}{\texttt{Player}}
\newcommand{\catalogservice}{\texttt{CatalogService}}
\newcommand{\playbackservice}{\texttt{PlaybackService}}
\newcommand{\libraryservice}{\texttt{LibraryService}}
\newcommand{\playliststore}{\texttt{PlaylistStore}}
\newcommand{\likedsongstore}{\texttt{LikedSongStore}}
\newcommand{\downscaler}{\texttt{downscaler}}
\newcommand{\lyricsparser}{\texttt{LyricsParser}}
\newcommand{\lyricsprovider}{\texttt{LyricsProvider}}
\newcommand{\streamresolver}{\texttt{StreamResolver}}
\newcommand{\ytmusicsource}{\texttt{YTMusicSource}}
\newcommand{\vlcaudiobackend}{\texttt{VLCAudioBackend}}
\newcommand{\nowplayingbar}{\texttt{NowPlayingBar}}
\newcommand{\lyricspanel}{\texttt{LyricsPanel}}
\newcommand{\albumart}{\texttt{AlbumArt}}
\newcommand{\tracklist}{\texttt{TrackList}}
\newcommand{\trackprogress}{\texttt{TrackProgress}}
\newcommand{\volumecontrol}{\texttt{VolumeControl}}
\newcommand{\playercontrols}{\texttt{PlayerControls}}
\newcommand{\sidebar}{\texttt{Sidebar}}
\newcommand{\searchbar}{\texttt{SearchBar}}
\newcommand{\queuepanel}{\texttt{QueuePanel}}
\newcommand{\nexttrackcard}{\texttt{NextTrackCard}}
\newcommand{\actionhints}{\texttt{ActionHints}}
\newcommand{\playbackcenter}{\texttt{PlaybackCenter}}
\newcommand{\playerutilities}{\texttt{PlayerUtilities}}
\newcommand{\progressrow}{\texttt{ProgressRow}}
\newcommand{\trackinfo}{\texttt{TrackInfo}}
\newcommand{\slider}{\texttt{Slider}}
\newcommand{\searchinput}{\texttt{SearchInput}}
\newcommand{\searchresult}{\texttt{SearchResults}}
\newcommand{\navitem}{\texttt{NavItem}}
\newcommand{\homesc}{\texttt{HomeScreen}}
\newcommand{\searchsc}{\texttt{SearchScreen}}
\newcommand{\librarysc}{\texttt{LibraryScreen}}
\newcommand{\playlistsc}{\texttt{PlaylistScreen}}
\newcommand{\fullscreensc}{\texttt{FullscreenScreen}}
\newcommand{\lyricssc}{\texttt{LyricsScreen}}

%% The macros above are kept verbatim.  As control words they swallow the space
%% that follows them ("\py 3.12" -> "Python3.12"), so each is wrapped with
%% \xspace *after* its original definition (definitions themselves untouched).
\usepackage{xspace}
\xspaceaddexceptions{(}
\newcommand{\xspacewrap}[1]{%
  \csletcs{orig@#1}{#1}%
  \csdef{#1}{\csuse{orig@#1}\xspace}}
\forcsvlist{\xspacewrap}{py,textual,vlc,sqlite,lrc,ansi,ytmusic,ytdlp,pil,tui,musictui,eventbus,protocol,track,lyric,playlist,playstate,repeatmode,audiobackend,catalogsource,trackcache,musicqueue,player,catalogservice,playbackservice,libraryservice,playliststore,likedsongstore,downscaler,lyricsparser,lyricsprovider,streamresolver,ytmusicsource,vlcaudiobackend,nowplayingbar,lyricspanel,albumart,tracklist,trackprogress,volumecontrol,playercontrols,sidebar,searchbar,queuepanel,nexttrackcard,actionhints,playbackcenter,playerutilities,progressrow,trackinfo,slider,searchinput,searchresult,navitem,homesc,searchsc,librarysc,playlistsc,fullscreensc,lyricssc}

%% Macros that the original text uses but never defined (added so it compiles)
\DeclareRobustCommand{\method}[1]{\texttt{\def\_{\textunderscore\allowbreak}#1}}
\newcommand{\property}[1]{\texttt{#1}}
\newcommand{\pyvlc}{\texttt{python-vlc}\xspace}
\newcommand{\numpy}{\texttt{NumPy}\xspace}
\newcommand{\ytmusicapi}{\texttt{ytmusicapi}\xspace}

%% ---------------------------------------------------------------------------
%% 18.  Title metadata
%% ---------------------------------------------------------------------------
\title{\textbf{\musictui: A Terminal-Based Music Player with Synced Lyrics,\\
ANSI Half-Block Colour Art, and Event-Driven Architecture}}
\author{\textbf{MusicTUI Project}\\[3pt]
  {\small Open-source software project}\\[1pt]
  {\small\href{https://github.com/RosarioLC/musictui}{\texttt{github.com/RosarioLC/musictui}}}}
\date{\today}

\begin{document}
\maketitle

\begin{abstract}
    This paper presents \musictui, a terminal-based music player built with
    \py 3.12+ and the \textual framework. The application streams music from
    YouTube Music, manages playlists and liked songs, displays synchronized
    lyrics, and renders album artwork as \ansi half-block colour text within
    a responsive terminal user interface. The system employs a clean layered
    architecture comprising presentation, service, domain, and storage layers,
    with \protocol-based dependency injection enabling extensibility and
    testability. An event-driven communication model using a central
    \eventbus decouples the playback engine from UI widgets, eliminating
    polling and reducing coupling. Domain models such as \track, \playstate,
    and \lyric are implemented as immutable frozen dataclasses with slot
    optimization for memory efficiency. The application leverages \pyvlc for
    audio playback, \sqlite for persistent storage with schema migration,
    \pil and \numpy for image downscaling to \ansi half-block colour art,
    and a triple-cached LRU cache for tracks, lyrics, and stream URLs.
    Synchronized lyrics are fetched via a multi-tier fallback system
    (\ytmusic \(\rightarrow\) lrclib \(\rightarrow\) syncedlyrics
    \(\rightarrow\) linear distribution). The paper documents the complete
    architecture, design decisions, implementation details, testing
    methodology, and future work for this open-source project.
\end{abstract}

\keywords{terminal user interface, music player, Python, Textual framework, synchronous lyrics, ANSI art, event-driven architecture, dependency injection, SQLite, VLC media library}

\section{Introduction}

Terminal-based applications have experienced a renaissance in recent years,
driven by developer familiarity with command-line environments, the ubiquity
of SSH-based remote workflows, and the growing capability of modern terminal
emulators to render rich colour graphics. Within this landscape, music players
have traditionally occupied a niche dominated by minimal, text-oriented tools
such as cmus, ncmpcpp, and mpc. While these tools are powerful, they often
lack the visual polish and interactive features that users expect from modern
media consumption software.

The \musictui project addresses this gap by delivering a feature-rich music
player that operates entirely within the terminal yet provides a visual
experience comparable to desktop applications. Built on the \textual framework
\cite{textual}, which provides a reactive \tui component model with CSS-like
styling, \musictui renders album artwork as full-colour \ansi half-block
characters, displays synchronized lyrics with auto-highlighting, and offers a
complete playlist management system -- all within a responsive, keyboard-driven
interface.

The primary contributions of this project are:

\begin{enumerate}
    \item A \protocol-based layered architecture that cleanly separates
          concerns across presentation, service, domain, and storage layers.
    \item An event-driven communication model using a central \eventbus that
          decouples the playback engine from UI components.
    \item A novel \ansi half-block colour image rendering pipeline using
          \pil and \numpy that converts album artwork into terminal-compatible
          colour art.
    \item A multi-tier lyrics fetching system with three fallback layers
          ensuring maximum coverage for synced lyrics.
    \item An immutable domain model using \py's frozen dataclasses with
          slot optimization for thread safety and memory efficiency.
    \item A comprehensive test suite covering all layers of the application
          with in-memory mocks and \sqlite databases.
\end{enumerate}

The remainder of this paper is organized as follows. Section~2 describes the
system architecture and design philosophy. Section~3 details the domain models.
Section~4 covers the playback subsystem. Section~5 discusses the catalog
integration with YouTube Music. Section~6 presents the UI architecture.
Section~7 documents the storage layer. Section~8 describes the testing
strategy. Section~9 covers installation and usage. Section~10 outlines future
work. Section~11 concludes.

\newpage

\section{System Architecture}

\musictui employs a four-layer architecture that enforces unidirectional data
flow and dependency inversion. Figure~\ref{fig:architecture} illustrates the
overall layered structure.

\begin{figure}[H]
    \centering
    \includegraphics[width=\textwidth]{architecture_overview}
    \caption{Layered architecture overview showing the four layers of \musictui
    and their interdependencies. Arrows indicate that higher layers depend on
    lower layers, with dependency injection flowing in the opposite direction.}
    \label{fig:architecture}
\end{figure}

\subsection{Layer Descriptions}

\paragraph{Presentation Layer (UI)}
The presentation layer comprises the \textual application class
\musictui(App), six screen classes (\homesc, \searchsc, \librarysc,
\playlistsc, \fullscreensc, \lyricssc), and twenty-two reusable widget
classes. This layer is responsible solely for rendering state and capturing
user input. It contains no business logic; all interaction with application
state is delegated to the service layer through the \eventbus.

The UI follows a grid-based layout with a two-column, three-row structure:
the left column contains a collapsible sidebar for navigation, while the right
column holds the main content area. The bottom row is occupied by the
persistent \nowplayingbar, which is composited across all screens. CSS-like
styling is defined in \texttt{style.tcss}, a 597-line theme file that
implements a dark Spotify-inspired palette with \#121212 background and
\#1ed760 green accent colour.

\paragraph{Service Layer}
The service layer consists of three service classes that aggregate domain
objects and coordinate cross-cutting concerns:

\begin{itemize}
    \item \catalogservice wraps a \catalogsource implementation and a
          \trackcache, providing search, track retrieval, lyrics fetching,
          and stream URL resolution with automatic caching.
    \item \playbackservice wraps the \player and \eventbus, exposing an
          event-driven API for playback control that screens and widgets
          use exclusively.
    \item \libraryservice wraps the \playliststore, providing playlist CRUD
          operations, liked-songs management, and song addition/removal.
\end{itemize}

\paragraph{Domain Layer}
The domain layer contains the core business models: \track, \lyric,
\playstate, \playlist, and \repeatmode. These models are immutable
(frozen dataclasses) and implement value-based equality and hashing. The
domain layer has no dependencies on external libraries or frameworks.

\paragraph{Storage Layer}
The storage layer provides persistent data management through \sqlite:
\playliststore handles playlist and song persistence with track resolution,
while \likedsongstore manages the liked-songs collection. The
\texttt{Database} class provides a context-manager-wrapped SQLite connection
with automatic schema creation and migration support.

\subsection{Dependency Injection}

\musictui uses \protocol-based dependency injection, a form of structural
typing where dependencies are specified as \py \protocol classes rather than
concrete implementations. Two key protocols define this pattern:

\begin{designbox}
\textbf{CatalogSource Protocol:} Any class implementing
\method{search(query)}, \method{get\_track(id)},
\method{get\_track\_stream\_url(track)}, and \method{get\_lyrics(track)}
can serve as a music catalog provider. The current implementation is
\ytmusicsource, but the protocol enables future providers (local file
system, Spotify, etc.) to be injected without modifying consuming code.
\end{designbox}

\begin{designbox}
\textbf{AudioBackend Protocol:} Any class implementing
\method{play(url)}, \method{pause()}, \method{resume()}, \method{stop()},
\method{seek(position\_ms)}, \method{get\_position()},
\method{get\_duration()}, \method{set\_volume(volume)},
\method{get\_volume()}, and properties \property{is\_playing} and
\property{is\_paused} can serve as the audio backend. The current
implementation is \vlcaudiobackend, but the protocol enables testing with
\texttt{FakeBackend} and future replacement with alternative backends.
\end{designbox}

The wiring diagram is as follows:

\begin{verbatim}
MusicTUI (App)
  ├── YTMusicSource ──┐
  ├── TrackCache      │──> CatalogService
  └───────────────────┘
  ├── VLCAudioBackend ──┐
  ├── MusicQueue        │──> Player ──> PlaybackService
  └─────────────────────┘
  ├── Database ──> PlaylistStore ──> LibraryService
  └── EventBus ──┘
\end{verbatim}

\section{Domain Models}

The domain models in \musictui are implemented as immutable frozen dataclasses
using \py 3.10+'s \texttt{@dataclass(slots=True, frozen=True)} decorator.
This design choice provides three key benefits:

\begin{enumerate}
    \item \textbf{Thread safety:} Frozen objects cannot be mutated after
          construction, eliminating race conditions in the multi-threaded
          \textual worker model.
    \item \textbf{Memory efficiency:} The \texttt{slots=True} parameter
          eliminates the \texttt{\_\_dict\_\_} per-instance dictionary,
          reducing memory overhead by approximately 40\% compared to
          standard dataclasses.
    \item \textbf{Functional updates:} The \method{replace(**changes)}
          pattern produces new instances with modified fields, following
          the immutable update idiom common in functional programming.
\end{enumerate}

\subsection{Track Model}

The \track dataclass is the central data structure of the application,
representing a single music track with full metadata:

\begin{lstlisting}[language=Python, caption={The \track frozen dataclass with
default value normalization and functional update support.}]
@dataclass(slots=True, frozen=True)
class Track:
    id: int
    title: str = "Unknown Title"
    artist: str = "Unknown Artist"
    album: str = "Unknown Album"
    duration_ms: int = 0
    image_path: str = "Unknown Image Path"
    lyrics_path: str = "Unknown Lyrics Path"
    url: str = ""
    provider_id: str = ""

    def __post_init__(self) -> None:
        if not self.title or not self.title.strip():
            object.__setattr__(self, "title", "Unknown Title")
        if self.duration_ms < 0:
            object.__setattr__(self, "duration_ms", 0)

    def to_dict(self) -> dict[str, object]:
        return {field: getattr(self, field)
                for field in self.__dataclass_fields__}

    def __hash__(self) -> int: return hash(self.id)
    def __eq__(self, other: object) -> bool:
        if not isinstance(other, Track): return NotImplemented
        return self.id == other.id

    def replace(self, **changes: object) -> Track:
        return replace(self, **changes)
\end{lstlisting}

Key design decisions:
\begin{itemize}
    \item Track IDs are computed as \texttt{zlib.crc32(video\_id) \& 0xFFFFFFFF},
          producing deterministic 32-bit unsigned integers from YouTube video IDs.
    \item The \method{to\_dict()} method enables straightforward serialization
          for database storage.
    \item Equality and hashing are based solely on \texttt{id}, allowing tracks
          to be used in sets and as dictionary keys.
\end{itemize}

\subsection{LyricLine Model}

The \lyric dataclass represents a single line of synchronized lyrics:

\begin{lstlisting}[language=Python, caption={The \lyric frozen dataclass.}]
@dataclass(slots=True, frozen=True)
class LyricLine:
    timestamp_ms: int
    text: str
\end{lstlisting}

\subsection{PlayState Model}

The \playstate dataclass captures a snapshot of playback state:

\begin{lstlisting}[language=Python, caption={The \playstate frozen dataclass.}]
@dataclass(slots=True, frozen=True)
class PlayState:
    is_playing: bool = False
    is_paused: bool = False

    def replace(self, **changes: object) -> PlayState:
        return replace(self, **changes)
\end{lstlisting}

The dual boolean fields allow all four states: not playing/playing,
paused/not paused. In practice, \texttt{is\_playing=True} and
\texttt{is\_paused=True} together represent the paused state.

\subsection{Playlist Model}

The \playlist class is a mutable ordered collection that maintains both a
list and a set of track IDs for O(1) membership testing:

\begin{archbox}
\textbf{Dual-structure invariant:} The \playlist class maintains
\property{songs} (a \texttt{list[Track]} for ordered iteration and
index-based access) and \property{song\_ids} (a \texttt{set[int]} for
O(1) membership checks). All mutations (\method{add\_song},
\method{remove\_song}) update both structures atomically, preserving
the invariant that \texttt{t.id in song\_ids} if and only if
\texttt{t in songs}.
\end{archbox}

The cursor-based navigation (\method{next\_song}, \method{previous\_song})
supports sequential track traversal within a playlist context.

\subsection{RepeatMode Enum}

The \repeatmode enumeration defines three repeat modes using modular arithmetic:

\begin{lstlisting}[language=Python, caption={RepeatMode enum with cycle logic.}]
class RepeatMode(Enum):
    OFF = 0
    ALL = 1
    ONE = 2

    def next_mode(self) -> RepeatMode:
        return RepeatMode((self.value + 1) % 3)
\end{lstlisting}

\begin{figure}[H]
    \centering
    \includegraphics[width=0.64\textwidth]{repeat_mode_diagram}
    \caption{The repeat mode cycle diagram showing the three states and
    transitions via the \method{next\_mode()} method.}
    \label{fig:repeat_mode}
\end{figure}

\section{Playback Subsystem}

The playback subsystem coordinates audio output, queue management, and
playback state. It consists of three core components: the audio backend
protocol, the music queue, and the high-level player controller.

\subsection{Audio Backend Protocol}

The \audiobackend \protocol defines the interface for all audio playback
backends. The current implementation, \vlcaudiobackend, wraps the
\pyvlc library \cite{vlc}:

\begin{lstlisting}[language=Python, caption={VLCAudioBackend wrapping python-vlc.}]
class VLCAudioBackend:
    def __init__(self) -> None:
        self._instance = vlc.Instance(
            "--quiet", "--no-video-title-show"
        )
        self._player = None

    def _ensure_player(self) -> Any:
        if self._player is None or not hasattr(self._player, "play"):
            self._player = self._instance.media_player_new()
        return self._player

    def play(self, url: str) -> None:
        player = self._ensure_player()
        media = self._instance.media_new(url)
        player.set_media(media)
        player.play()

    def seek(self, position_ms: int) -> None:
        if self._player:
            self._player.set_time(position_ms)

    @property
    def is_playing(self) -> bool:
        val = self._player.is_playing()
        return bool(val) if val is not None else False
\end{lstlisting}

The backend uses lazy initialization of the \texttt{MediaPlayer} instance
and caches the \texttt{Paused} state enum for fast comparison in the
\property{is\_paused} property.

\subsection{Music Queue}

The \musicqueue manages the playback queue with history tracking, repeat
modes, and shuffle support:

\begin{archbox}
\textbf{Queue invariants:} The queue maintains three internal structures:
\begin{itemize}
    \item \texttt{\_tracks}: upcoming tracks to play (list)
    \item \texttt{\_track\_ids}: track IDs for O(1) lookups (set)
    \item \texttt{\_history}: played tracks for back navigation (list)
    \item \texttt{\_current}: the currently playing track
\end{itemize}
The \method{play\_next()} method transfers the current track to history
and pops the next track from the queue. In \repeatmode.ALL, when the
queue is exhausted, the history is restored to the queue for a full repeat.
\end{archbox}

\begin{lstlisting}[language=Python, caption={Core queue operations.}]
class MusicQueue:
    def play_next(self) -> Track | None:
        if not self._tracks:
            if self._repeat == RepeatMode.ALL and self._history:
                self._tracks = self._history
                self._history = []
            else:
                return None
        if self._current is not None:
            self._history.append(self._current)
        self._current = self._tracks.pop(0)
        self._track_ids.discard(self._current.id)
        return self._current

    def skip_to(self, track: Track) -> None:
        if track.id not in self._track_ids:
            raise ValueError("Track is not in the queue")
        if self._current is not None:
            self._history.append(self._current)
        self._tracks.remove(track)
        self._track_ids.discard(track.id)
        self._current = track
\end{lstlisting}

\subsection{Player Controller}

The \player class sits at the intersection of the audio backend, the queue,
and the UI. It handles track transitions, repeat logic, and UI event
emission:

\begin{lstlisting}[language=Python, caption={Player controller coordinating
backend, queue, and UI events.}]
class Player:
    def play(self, track: Track) -> None:
        if not self._queue.contains(track):
            self._queue.add(track)
        self._queue.skip_to(track)
        url = (self._stream_resolver(track)
               if self._stream_resolver else track.url)
        self._backend.play(url)
        self._schedule_ui_update()

    def _update_ui(self) -> None:
        if self._app is None: return
        current = self.current_track
        if current is None: return
        self._app.post_message(
            PlaybackStateChanged(self.is_playing, self.is_paused)
        )
        self._app.post_message(
            PositionUpdated(self.position_ms, self.duration_ms)
        )
        self._app.post_message(QueueUpdated(self.queue))
\end{lstlisting}

The \method{on\_track\_ended()} method implements the repeat logic:
\repeatmode.ONE restarts the current track, while \repeatmode.ALL and
\repeatmode.OFF advance to the next track in the queue.

\newpage

\section{Catalog Integration}

The catalog subsystem provides music discovery and metadata retrieval through
integration with YouTube Music and supplementary lyrics APIs.

\subsection{YouTube Music Source}

The \ytmusicsource class implements the \catalogsource \protocol using the
\ytmusicapi library \cite{ytmusicapi}:

\begin{lstlisting}[language=Python, caption={YTMusicSource search and track construction.}]
class YTMusicSource(CatalogSource):
    def search(self, query: str) -> list[Track]:
        client = self._get_client()
        if client is None: return []
        results = client.search(query, filter="songs")
        tracks: list[Track] = []
        seen: set[str] = set()
        for result in results:
            if not result.get("videoId"): continue
            vid = result["videoId"]
            if vid in seen: continue
            seen.add(vid)
            tracks.append(self._make_track(result))
            if len(tracks) >= 12: break
        return tracks

    @staticmethod
    def _make_track(result: dict) -> Track:
        vid = result["videoId"]
        return Track(
            id=zlib.crc32(vid.encode("utf-8")) & 0xFFFFFFFF,
            provider_id=vid,
            title=result.get("title", "Unknown Title"),
            artist=artists[0]["name"] if artists else "Unknown Artist",
            album=(result.get("album") or {}).get("name", "Unknown Album"),
            duration_ms=_duration_to_ms(result.get("duration")),
            image_path=thumbs[-1]["url"] if thumbs else "",
            url=f"https://music.youtube.com/watch?v={vid}",
        )
\end{lstlisting}

\newpage

\subsection{Lyrics Fetching Pipeline}

The lyrics fetching system implements a four-tier fallback pipeline to
maximize coverage:

\begin{figure}[H]
    \centering
    \includegraphics[width=\textwidth]{lyrics_pipeline}
    \caption{Lyrics fetching pipeline: a four-tier fallback system that maximises
    coverage by attempting YouTube Music, then lrclib.net, then syncedlyrics,
    and finally falling back to linear time distribution.}
    \label{fig:lyrics_pipeline}
\end{figure}

\paragraph{Tier 1: YouTube Music (ytmusicapi)}
The system first attempts to retrieve lyrics via the YouTube Music watch
playlist endpoint. It fetches the watch playlist for the track's video ID,
extracts the lyrics bid (identifier), and then retrieves the full lyrics
text.

\paragraph{Tier 2: lrclib.net API}
When YouTube Music lyrics are unavailable, the system queries the lrclib
API \cite{lrclib} at \texttt{https://lrclib.net/api/search} with the
artist and track name. The response (JSON containing synced or plain
lyrics) is parsed using the \lyricsparser.

\paragraph{Tier 3: syncedlyrics Library}
As a third fallback, the \texttt{syncedlyrics} library is invoked to search
for synced lyrics using the artist and track name as a query.

\paragraph{Tier 4: Linear Distribution}
When all API-based sources fail, plain lyrics text is distributed evenly
across the track duration using linear interpolation:

\begin{equation}
    \text{timestamp}_i = \left\lfloor \frac{i}{N-1} \times \text{duration\_ms} \right\rfloor
\end{equation}

where $N$ is the number of lyric lines and $i$ ranges from $0$ to $N-1$.

\subsection{LRC Parsing}

The \lyricsparser converts raw \lrc-formatted text into a list of \lyric
objects using regex matching:

\begin{lstlisting}[language=Python, caption={LRC timestamp parsing.}]
_LRC_PATTERN = re.compile(
    r"^\[(\d{1,2}):(\d{2})(?:\.(\d{1,3}))?\]\s*(.*)$"
)

@staticmethod
def parse(raw_text: str) -> list[LyricLine]:
    lines: list[LyricLine] = []
    for raw_line in raw_text.splitlines():
        match = _LRC_PATTERN.match(raw_line.strip())
        if match is None: continue
        minutes, seconds, fraction, text = match.groups()
        milliseconds = int((fraction or "0").ljust(3, "0"))
        timestamp_ms = (int(minutes) * 60_000
                       + int(seconds) * 1_000
                       + milliseconds)
        lines.append(LyricLine(timestamp_ms, text))
    return lines
\end{lstlisting}

The regex handles optional centiseconds (the \texttt{.xx} portion),
right-padding them with zeros to form a three-digit millisecond value.

\subsection{Stream Resolution}

Stream URLs are resolved via lazy-loaded \ytdlp \cite{ytdlp}. The
\texttt{yt-dlp} module is imported only when needed, and the result is
cached in the \trackcache keyed by the track's URL:

\begin{lstlisting}[language=Python, caption={Lazy yt-dlp stream resolution.}]
def _get_yt_dlp_dl(self) -> Any:
    if self._yt_dlp_dl is not None:
        return self._yt_dlp_dl
    try:
        yt_dlp = import_module("yt_dlp")
        options = {
            "format": "bestaudio/best",
            "quiet": True,
            "noplaylist": True,
            "cachedir": True,
        }
        self._yt_dlp_dl = yt_dlp.YoutubeDL(options)
    except ImportError:
        self._yt_dlp_dl = None
    return self._yt_dlp_dl

def get_track_stream_url(self, track: Track) -> str:
    dl = self._get_yt_dlp_dl()
    if dl is None:
        return track.url
    info = dl.extract_info(track.url, download=False)
    return info.get("url", track.url)
\end{lstlisting}

\section{User Interface}

The UI is built on the \textual framework, which provides a reactive component
model with CSS-like styling, keyboard bindings, and built-in support for
asynchronous workers.

\subsection{Application Structure}

The \musictui(App) class manages the application lifecycle, service
initialization, screen navigation, and asset preloading:

\begin{lstlisting}[language=Python, caption={Main application class wiring
all services together.}]
class MusicTUI(App):
    CSS_PATH = "style.tcss"
    BINDINGS = [
        Binding("q", "quit", "Quit"),
        Binding("h", "go_home", "Home"),
        Binding("s", "go_search", "Search"),
        Binding("l", "go_library", "Library"),
        Binding("f", "toggle_fullscreen", "Fullscreen"),
    ]

    def on_mount(self) -> None:
        self._event_bus = EventBus(self)
        self._catalog_source = YTMusicSource()
        self._store = Database("store.db")
        self._track_cache = TrackCache(max_size=50)
        self._playlist_store = PlaylistStore(
            self._store, self._catalog_source
        )
        self._catalog_service = CatalogService(
            self._catalog_source, self._track_cache
        )
        self._library_service = LibraryService(self._playlist_store)
        self._player = Player(
            backend=VLCAudioBackend(),
            queue=MusicQueue(),
            stream_resolver=lambda t:
                self._catalog_service.get_stream(t),
        )
        self._player.set_textual_app(self)
        self._playback_service = PlaybackService(
            self._player, self._event_bus
        )
        self.push_screen(HomeScreen(
            self._catalog_service,
            self._playback_service,
            self._library_service
        ))
\end{lstlisting}

\subsection{Screen Architecture}

\musictui defines six screen classes, each responsible for a specific
view of the application:

\begin{table}[H]
\centering
\small
\caption{Screen overview table describing each screen's purpose and key
features.}
\label{tab:screens}
\begin{tabularx}{\textwidth}{lX}
\toprule
\textbf{Screen} & \textbf{Description} \\
\midrule
\homesc & Recommended tracks with infinite scroll, search bar,
next-track card, and now-playing bar \\
\searchsc & Query input with search results in a track list, supporting
like and add-to-playlist actions \\
\librarysc & Two-panel layout: playlist table (left) and song list (right)
with CRUD operations \\
\playlistsc & Single playlist view with track playback, navigation, and
song removal \\
\fullscreensc & Large album-art overlay filling the terminal, with track
title/artist caption \\
\lyricssc & Synced lyrics overlay with auto-highlighting, loaded via
background worker \\
\bottomrule
\end{tabularx}

\end{table}

\begin{figure}[H]
    \centering
    \includegraphics[width=0.9\textwidth, height=6.5cm, keepaspectratio]{ui_layout.png}
    \caption{The home screen layout showing the two-column grid with sidebar
navigation, recommended tracks list, next-track preview, and the persistent
now-playing bar at the bottom.}
    \label{fig:ui_layout}
\end{figure}

\subsection{Widget Hierarchy}

The application uses 22 reusable widget classes. Figure~\ref{fig:ansi_art_example}
shows the main groupings.

\begin{figure}[H]
    \centering
    \includegraphics[width=\textwidth]{ansi_art_example}
    \caption{ANSI half-block colour rendering: each terminal character
combines two pixel rows -- the top row as background colour and the bottom
row as foreground colour -- producing a dense, full-colour image.}
    \label{fig:ansi_art_example}
\end{figure}

\subsubsection{Now-Playing Bar}

The \nowplayingbar is the most complex widget, composed of three sub-sections:

\begin{itemize}
    \item \textbf{Track info} (\trackinfo): Album art (\albumart) +
          title/artist labels. Reacts to \eventbus messages to update
          displayed information.
    \item \textbf{Playback center} (\playbackcenter): Transport controls
          (\playercontrols: prev, play/pause, next) + seek bar
          (\progressrow: current time + \trackprogress slider + duration).
    \item \textbf{Utilities} (\playerutilities): Lyrics toggle button +
          \volumecontrol (mute button + volume slider) + fullscreen toggle.
\end{itemize}

The bar uses \textual's reactive properties to respond to state changes.

\subsubsection{Album Art Rendering}

The \albumart widget renders album artwork as \ansi half-block colour text.
It uses a background worker to perform the computationally intensive
downscaling operation.

\begin{lstlisting}[language=Python, caption={AlbumArt widget rendering.}]
class AlbumArt(Static):
    def _render_artwork(self) -> None:
        self._render_generation += 1
        cols = self._art_cols or round(
            self.size.height * CELL_ASPECT
        )
        self.run_worker(
            partial(downscale, self._source,
                    cols=cols, cell_aspect=CELL_ASPECT),
            name="album-art-render", thread=True,
        )

    def on_worker_state_changed(self, event):
        if event.state == WorkerState.SUCCESS:
            self.update(Text.from_ansi(event.worker.result))
\end{lstlisting}

\subsection{ANSI Half-Block Colour Rendering}

The \downscaler module converts any image (local file or URL) to \ansi
half-block colour text using \pil and \numpy \cite{numpy,pillow}. The algorithm:

\begin{enumerate}
    \item Load the image using \pil (Pillow).
    \item Resize to $W$ columns by $H = 2 \times \text{rows}$ pixels
          using BICUBIC interpolation.
    \item For each pair of pixel rows $(2i, 2i+1)$, combine the top pixel
          (background colour) and bottom pixel (foreground colour) into a
          single \ansi half-block character.
    \item Each character is encoded as an escape sequence:
          \texttt{\textbackslash{}e[48;2;R;G;Bm\textbackslash{}e[38;2;r;g;bm\textbackslash{}2584}.
\end{enumerate}

\begin{lstlisting}[language=Python, caption={The core downscaling algorithm.}]
_ANSI_FMT = "\x1b[48;2;{};{};{}m\x1b[38;2;{};{};{}m\u2584"

def downscale(image_path: str, cols: int = 24,
              cell_aspect: float = 2.1) -> str:
    text_rows = max(1, round(cols / cell_aspect))
    W, H = cols, text_rows * 2
    image = Image.open(BytesIO(_read_image(image_path)))
    image = image.convert("RGB")
    image = image.resize((W, H), Image.BICUBIC)
    arr = np.asarray(image)
    lines = []
    for ty in range(0, H, 2):
        top = arr[ty].tolist()
        bot = arr[ty + 1].tolist()
        line_parts = [
            _ANSI_FMT.format(a[0], a[1], a[2], b[0], b[1], b[2])
            for a, b in zip(top, bot)
        ]
        lines.append("".join(line_parts))
    return "\n".join(lines)
\end{lstlisting}

The algorithm is cached via \texttt{@lru\_cache(maxsize=64)}, avoiding
redundant computation.

\begin{figure}[H]
    \centering
    \includegraphics[width=0.9\textwidth]{lyrics_screen.png}
    \caption{The lyrics overlay screen showing synchronized lyrics with the
    currently active line highlighted in green, loaded via a background
    worker and highlighted every 250ms using binary search.}
    \label{fig:lyrics_screen}
\end{figure}

\subsection{Event-Driven Communication}

The \eventbus decouples the player from UI widgets. Instead of widgets
polling the player for state, the player posts messages to the \eventbus,
and widgets subscribe to the relevant message types:

\begin{lstlisting}[language=Python, caption={Event types and EventBus.}]
class EventBus:
    def __init__(self, app: App) -> None:
        self._app = app

    def emit(self, message: Message) -> None:
        self._app.post_message(message)

    def emit_playback_started(self, track: Track) -> None:
        self.emit(PlaybackStarted(track))

    def emit_position_updated(
        self, position_ms: int = 0, duration_ms: int = 0
    ) -> None:
        self.emit(PositionUpdated(position_ms, duration_ms))
\end{lstlisting}

Message types include: \texttt{PlaybackStarted}, \texttt{PlaybackEnded},
\texttt{PlaybackStateChanged}, \texttt{PositionUpdated}, \texttt{QueueUpdated},
\texttt{LikedSongsChanged}, and \texttt{PlaylistChanged}.

\begin{figure}[H]
    \centering
    \includegraphics[width=\textwidth]{event_flow}
    \caption{Sequence diagram showing the event flow when a user clicks a
    track: the click propagates through the track list widget to the screen,
    which calls the playback service, which notifies the player, which posts
    messages to the event bus, which are consumed by all subscribed widgets.}
    \label{fig:event_flow}
\end{figure}

\newpage

\subsection{Library Screen}

The \librarysc provides a two-panel layout for playlist management. The left
panel is a \texttt{DataTable} listing all playlists, while the right panel
displays the songs in the selected playlist:

\begin{figure}[H]
    \centering
    \includegraphics[width=0.9\textwidth]{library_screen.png}
    \caption{The library screen with the two-panel layout: playlist table on
    the left and song list on the right. Supports playlist CRUD operations
    (create, rename, delete) via modal dialogs.}
    \label{fig:library_screen}
\end{figure}

\subsection{Modal Dialogs}

Three modal screen classes handle user input: \texttt{NewPlaylistScreen},
\texttt{RenamePlaylistScreen}, and \texttt{PlaylistPickerScreen}. All modals
support \texttt{Escape}-to-cancel and \texttt{Enter}-to-submit shortcuts.

\subsubsection{Fullscreen Screen}

The \fullscreensc provides a full-terminal overlay showing a large
\ansi half-block colour rendering of the current track's album art:

\begin{figure}[H]
    \centering
    \includegraphics[width=0.9\textwidth]{fullscreen_art.png}
    \caption{The fullscreen overlay displaying a large \ansi half-block colour
    rendering of the current track's album art, filling the terminal window.}
    \label{fig:fullscreen_art}
\end{figure}

\section{Storage Layer}

The storage layer provides persistent data management using \sqlite \cite{sqlite},
a lightweight, serverless, zero-configuration SQL database engine.

\subsection{Database Wrapper}

The \texttt{Database} class wraps an \sqlite connection with context manager
support and automatic schema management:

\begin{lstlisting}[language=Python, caption={Database wrapper with migration.}]
class Database:
    def __init__(self, db_path: str) -> None:
        self._db_path = Path(db_path)
        self._connect()

    def _connect(self) -> None:
        self._conn = sqlite3.connect(str(self._db_path))
        self._conn.row_factory = sqlite3.Row
        rows = self._conn.execute(
            "PRAGMA table_info(liked_songs)"
        ).fetchall()
        if "track_id" in [r["name"] for r in rows]:
            self._conn.execute("DROP TABLE IF EXISTS liked_songs")
        self._conn.execute(
            "CREATE TABLE IF NOT EXISTS liked_songs "
            "(provider_id TEXT PRIMARY KEY)"
        )
        self._conn.execute(
            "CREATE TABLE IF NOT EXISTS tracks "
            "(id INTEGER PRIMARY KEY AUTOINCREMENT, "
            "provider_id TEXT UNIQUE, title TEXT, artist TEXT, "
            "album TEXT, duration_ms INTEGER, "
            "image_path TEXT, lyrics_path TEXT, url TEXT)"
        )
        self._conn.execute(
            "CREATE TABLE IF NOT EXISTS playlists "
            "(id INTEGER PRIMARY KEY AUTOINCREMENT, name TEXT NOT NULL)"
        )
        self._conn.execute(
            "CREATE TABLE IF NOT EXISTS playlist_songs "
            "(playlist_id INTEGER NOT NULL, track_id INTEGER NOT NULL, "
            "FOREIGN KEY(playlist_id) REFERENCES playlists(id), "
            "FOREIGN KEY(track_id) REFERENCES tracks(id), "
            "PRIMARY KEY(playlist_id, track_id))"
        )
        self._conn.commit()

    def __enter__(self) -> Self: return self
    def __exit__(self, *args: object) -> None: self.close()
\end{lstlisting}

The schema includes four tables: \texttt{tracks}, \texttt{playlists},
\texttt{playlist\_songs} (junction table with foreign keys), and
\texttt{liked\_songs}.

\subsection{Playlist and Liked-Songs Stores}

The \playliststore class manages playlist and song persistence, resolving
stored track IDs back to \track objects. The \likedsongstore manages the
liked-songs collection using \texttt{INSERT OR IGNORE} to prevent duplicates.

\section{Testing}

\musictui includes a comprehensive test suite organized in parallel to the
production code structure, using \texttt{pytest} \cite{pytest}.

\subsection{Test Structure}

\begin{verbatim}
tests/
├── catalog/
│   ├── fake_catalog_source.py      # In-memory mock for CatalogSource
│   ├── test_fake_catalog_source.py # Search functionality tests
│   ├── test_lyrics_parser.py       # LRC parsing edge cases
│   ├── test_lyrics_provider.py     # Lyrics fetch tests
│   └── test_track_cache.py         # LRU cache + TrackCache tests
├── models/
│   ├── test_play_state.py          # Immutable state tests
│   ├── test_playlist.py            # CRUD + navigation tests
│   ├── test_repeat_mode.py         # Cycle tests
│   └── test_track.py               # Construction + replace tests
├── playback/
│   ├── fake_backend.py             # In-memory mock for AudioBackend
│   ├── test_music_queue.py         # Queue operation tests
│   └── test_player.py              # Player operation tests
└── store/
    ├── test_liked_songs_store.py   # SQLite in-memory tests
    └── test_playlist_store.py      # Playlist persistence tests
\end{verbatim}

\subsection{Mock Implementations}

Two key mock classes enable unit testing without external dependencies:

\begin{lstlisting}[language=Python, caption={FakeBackend mock for testing.}]
class FakeBackend(AudioBackend):
    def __init__(self):
        self._playing = False
        self._paused = False
        self._volume = 75
        self._position_ms = 0

    def play(self, url: str) -> None:
        self._playing = True
        self._paused = False

    def pause(self) -> None: self._paused = True
    def resume(self) -> None: self._paused = False
    def stop(self) -> None: self._playing = False

    def seek(self, position_ms: int) -> None:
        self._position_ms = position_ms

    def get_position(self) -> int: return self._position_ms
    def get_duration(self) -> int: return 180000
\end{lstlisting}

\subsection{Test Categories}

\paragraph{Model Tests (36 tests)}
Model tests verify frozen dataclass construction, the \method{replace()}
functional update pattern, immutability guarantees, hash/equality semantics,
and normalization edge cases. The \track model is tested for all default
values, field normalization (whitespace-only fields fall back to defaults),
negative duration clamping, CRC32 ID generation, and hash consistency for use
in sets and dictionaries. The \playlist model is tested for O(1) membership
via set-backed lookup, cursor position tracking on add/remove, and the
\method{next\_song} / \method{previous\_song} navigation methods.

\begin{lstlisting}[language=Python, caption={Immutability test for frozen dataclass.}]
class TestTrack:
    def test_immutability(self):
        t = Track(id=1, title="Test")
        try:
            t.title = "Changed"
            assert False, "Should not mutate frozen dataclass"
        except Exception:
            pass

    def test_hash_and_equality(self):
        t1 = Track(id=1, title="Test")
        t2 = Track(id=1, title="Different")
        assert t1 == t2  # Same id => equal
        assert hash(t1) == hash(t2)
        s = {t1}
        assert t2 in s  # Works in sets
\end{lstlisting}

\paragraph{Cache Tests (16 tests)}
The \trackcache and its underlying \texttt{\_LRUCache} generic LRU cache are
tested for eviction policy, LRU promotion (accessing a key moves it to the
most-recently-used position), key normalization for lyrics (whitespace
stripping and case-insensitive matching), and independent sub-caches for
tracks, lyrics, and stream URLs.

\begin{lstlisting}[language=Python, caption={LRU promotion test.}]
class TestLRUCache:
    def test_lru_promotion(self):
        cache = _LRUCache(max_size=3)
        cache.put("a", 1)
        cache.put("b", 2)
        cache.put("c", 3)
        cache.get("a")  # promotes "a" to MRU
        cache.put("d", 4)  # evicts "b" (was LRU)
        assert cache.get("a") == 1
        assert cache.get("b") is None  # evicted
\end{lstlisting}

\paragraph{Queue and Player Tests (41 tests)}
Queue tests verify FIFO ordering, history management, repeat mode cycling,
\method{skip\_to()} validation, duplicate prevention, shuffle, and the
\method{replace\_upcoming()} operation. Player tests use \texttt{FakeBackend}
to verify play/pause toggling, seek, volume control, position tracking,
stream resolver integration, and the \method{on\_track\_ended()} repeat
handling logic.

\paragraph{Catalog Mock Tests (8 tests)}
The \texttt{FakeCatalogSource} mock is tested for case-insensitive search
across title, artist, and album fields, empty query handling, and multiple
match returns.

\paragraph{Storage Tests (24 tests)}
Storage tests use in-memory \sqlite databases (\texttt{Database(":memory:")})
to verify CRUD operations, foreign key constraints, provider-based track
resolution (ensuring \texttt{zlib.crc32} in-memory IDs resolve correctly to
the database primary key), and the liked-songs store ordering.

\subsection{Running Tests}

\begin{lstlisting}[style=console, caption={Running the test suite.}]
$ cd music/
$ python -m pytest tests/ -v
======================== 167 passed in 0.16s =========================
$ python -m pytest tests/ -v --cov=music --cov-report=term-missing
\end{lstlisting}

The test suite covers all production modules in parallel: catalog (parsing,
caching, search mocks), models (frozen dataclasses, enum cycling), playback
(queue operations, player controls), and store (SQLite persistence). Tests
run in under one second, enabling rapid feedback during development.

\section{Getting Started}

\subsection{Prerequisites}

\begin{itemize}
    \item \py 3.12 or later
    \item \vlc media library (system package)
    \item A terminal emulator with 24-bit colour support
          (e.g., Windows Terminal, iTerm2, GNOME Terminal)
\end{itemize}

\subsection{Installation}

\begin{lstlisting}[style=console, caption={Installation instructions.}]
$ git clone https://github.com/username/musictui.git
$ cd music
$ python -m venv .venv
$ source .venv/bin/activate  # On Windows: .venv\Scripts\activate
$ pip install -e ".[dev]"
$ python -m pytest tests/ -v
\end{lstlisting}

\subsection{Running the Application}

\begin{lstlisting}[style=console, caption={Launching the application.}]
$ python main.py
\end{lstlisting}

\subsection{Keyboard Shortcuts}

\begin{table}[H]
\centering
\caption{Keyboard shortcuts available in \musictui.}
\label{tab:keybindings}
\begin{tabular}{ccp{8cm}}
\toprule
\textbf{Key} & \textbf{Mode} & \textbf{Action} \\
\midrule
\texttt{h} & Always & Navigate to Home screen \\
\texttt{s} & Always & Navigate to Search screen \\
\texttt{l} & Always & Navigate to Library screen \\
\texttt{f} & Always & Toggle fullscreen album art \\
\texttt{q} & Always & Quit application \\
\texttt{a} & Home/Search & Add currently playing track to playlist \\
\texttt{n} & Library & Create new playlist \\
\texttt{d} & Library/Playlist & Delete playlist / Remove song \\
\texttt{r} & Library & Rename selected playlist \\
\texttt{Escape} & Lyrics/Fullscreen & Close overlay screen \\
\texttt{Enter} & Search bar & Submit search query \\
\texttt{Escape} & Search input & Clear search input \\
\bottomrule
\end{tabular}

\end{table}

\section{Future Work}

Several areas have been identified for future enhancement of \musictui:

\subsection{Additional Music Providers}
The \catalogsource \protocol is designed to support multiple backends.
Planned implementations include: local file system provider with
\texttt{mutagen} metadata extraction, Spotify integration via \texttt{spotipy},
and Deezer integration via \texttt{deezer-python}.

\subsection{Enhanced Lyrics Features}
\begin{itemize}
    \item \textbf{Karaoke-style word-level highlighting:} Current lyrics
          display line-level highlighting. Word-level timing data from
          lrclib's \texttt{syncedLyricsPlain} field could enable per-word
          highlighting.
    \item \textbf{Offline lyrics caching:} Cache fetched lyrics to disk
          for offline playback.
    \item \textbf{Multi-language lyrics:} Display lyrics in multiple
          languages simultaneously when available.
\end{itemize}

\subsection{Audio Enhancements}
\begin{itemize}
    \item \textbf{Equalizer:} Real-time audio equalization using \vlc's
          built-in equalizer presets.
    \item \textbf{Crossfade:} Smooth transitions between tracks.
    \item \textbf{Audio visualization:} Real-time waveform or frequency
          spectrum visualization using \numpy's FFT capabilities.
    \item \textbf{Multi-backend support:} Add \texttt{pygame.mixer} or
          \texttt{portaudio} backends for systems where \vlc is unavailable.
\end{itemize}

\subsection{UI Improvements}
\begin{itemize}
    \item \textbf{Theme system:} Allow users to define custom colour themes
          via a \texttt{themes/} directory of \texttt{.tcss} files.
    \item \textbf{Responsive layouts:} Adapt the grid layout for small
          terminals (below 80 columns) by hiding the sidebar.
    \item \textbf{Mouse support:} Full mouse interaction for track selection,
          slider dragging, and button clicking across all screens.
    \item \textbf{Recently played screen:} Display a history of recently
          played tracks with the ability to re-play.
\end{itemize}

\subsection{Data and Analytics}
\begin{itemize}
    \item \textbf{Listening statistics:} Track total listening time, most
          played artists, and most played tracks.
    \item \textbf{Export/import:} Export playlists as \texttt{.m3u} or
          \texttt{.xspf} files; import from \texttt{.m3u}, \texttt{.xspf},
          or Spotify CSV exports.
\end{itemize}

\section{Conclusion}

This paper has presented \musictui, a terminal-based music player that
demonstrates how modern Python frameworks can produce rich, interactive
applications within the terminal environment. The project showcases several
key software engineering practices:

\begin{itemize}
    \item \textbf{Protocol-based architecture} provides a clean separation
          of concerns while maintaining flexibility for future extensions.
    \item \textbf{Event-driven communication} through the \eventbus eliminates
          tight coupling between components and enables independent evolution
          of the playback engine and UI layers.
    \item \textbf{Immutable domain models} using frozen dataclasses with
          slots ensure thread safety and memory efficiency in a
          multi-threaded environment.
    \item \textbf{ANSI half-block colour art} provides a unique visual
          experience that leverages the terminal's capabilities in a way
          that desktop GUI applications cannot replicate.
    \item \textbf{Multi-tier lyrics fetching} ensures maximum coverage for
          synchronized lyrics through a systematic fallback chain.
    \item \textbf{Comprehensive testing} with 167 unit tests across 10 test
          modules, using in-memory mocks and SQLite in-memory databases,
          provides confidence in the correctness of the application logic
          without requiring external dependencies.
\end{itemize}

The \musictui project demonstrates that terminal-based applications can
compete with desktop GUI applications in terms of features and user
experience, while maintaining the advantages of terminal applications:
lightweight resource usage, SSH accessibility, and keyboard-driven
efficiency.

\begin{thebibliography}{99}

\bibitem{textual}
Textual Framework.
\newblock \texttt{https://textual.textualize.io/}
\newblock A library for building terminal user interfaces in Python.

\bibitem{vlc}
VideoLAN.
\newblock \texttt{https://www.videolan.org/vlc/}
\newblock VLC media library -- a free and open-source cross-platform multimedia
framework and player.

\bibitem{sqlite}
SQLite Consortium.
\newblock \texttt{https://www.sqlite.org/}
\newblock SQLite -- A C-language library that implements a small, fast,
self-contained, high-reliability, full-featured, SQL database engine.

\bibitem{pillow}
Pillow (PIL Fork).
\newblock \texttt{https://pillow.readthedocs.io/}
\newblock The Python Imaging Library -- adds image processing capabilities to
your Python interpreter.

\bibitem{numpy}
Harris, C. R., Millman, K. J., van der Walt, S. J., Gommers, R., Virtanen, P.,
Cournapeau, D., \& Oliphant, T. E.
\newblock Array programming with \texttt{numpy}.
\newblock \emph{Nature}, 585(7825), 357--362, 2020.

\bibitem{ytmusicapi}
nathom.
\newblock \texttt{https://github.com/nathom/ytmusicapi}
\newblock A Python API for interacting with YouTube Music.

\bibitem{ytdlp}
yt-dlp contributors.
\newblock \texttt{https://github.com/yt-dlp/yt-dlp}
\newblock A youtube-dl fork with additional features and fixes.

\bibitem{lrclib}
lrclib.
\newblock \texttt{https://lrclib.net/}
\newblock A community-driven API for obtaining synced and plain lyrics.

\bibitem{pytest}
pytest contributors.
\newblock \texttt{https://docs.pytest.org/}
\newblock A framework for writing simple and scalable tests in Python.

\bibitem{pyvlc}
Pierre Raybaut.
\newblock \texttt{https://github.com/nickolag/python-vlc}
\newblock Python bindings for the VLC media library.

\bibitem{python}
Python Software Foundation.
\newblock \texttt{https://www.python.org/}
\newblock Python 3.12 -- What's New in Python 3.12.

\end{thebibliography}

\end{document}
